\section{Theory}
\label{sec:theory}

This section asks when the anchored cell model can recover the latent set $A_c$ from readings, and how many readings it needs. All statements are per cell, so we drop the index $c$. Proofs are in \cref{app:proofs}.

\subsection{Assumptions}

\begin{description}[leftmargin=2.2em,style=nextline]
\item[(A1) Independence.] Given $A$ and $\Theta$, the readings of one model are independent. They need not be identically distributed across Tell variants. $P_m$ then denotes the variant-averaged reading distribution of model $m$.
\item[(A2) Accuracy parametrisation.] If $A\neq\emptyset$, then $P_m(A)=p_m$, faithful readings are uniform on $A$, and $|A|\le2$.
\item[(A3) Shared errors.] The error distribution $\varepsilon$, restricted to $V\setminus A$, is the same for every model at the cell.
\item[(A4) Monotone accuracy.] $p_{m,c}=\sigma(\alpha_m-\beta_{g(c)})$. The order of models by accuracy is therefore the same at every cell.
\item[(A5) Silence.] If $A=\emptyset$, every model reads from $\varepsilon$: competence does not help where $\Tell$ says nothing.
\end{description}

(A2) and (A4) are modelling choices. (A1) holds approximately because readings come from separate runs with distinct seeds and variants. (A3) and (A5) are the substantive assumptions, and both may fail. Errors of different LLMs are known to be correlated~\cite{kim2025correlated,goel2025great}, which favours (A3) at the population level, but the \emph{shape} of errors is not flat. A strong model may also guess a value that $\Tell$ does not state, from domain knowledge. We return to both failures in \cref{sec:theory-scope} and test them in simulation.

\subsection{Accuracy-gradient signatures}

The key observation is that each defect form changes the reading distribution in a different way as accuracy grows. Take two models with $p_1>p_2$, write $\Delta=p_1-p_2$, and let
\begin{equation}
\label{eq:contrast}
D(y)=P_1(y)-P_2(y).
\end{equation}

\begin{lemma}[Signatures]
\label{lem:signatures}
Under (A2), (A3) and (A5), let $\tilde\varepsilon=\varepsilon(\cdot\mid V\setminus A)$. Then
\begin{equation}
\label{eq:signature}
D(y)=
\begin{cases}
\Delta\Bigl(\dfrac{\one[y\in A]}{|A|}-\tilde\varepsilon(y)\Bigr), & A\neq\emptyset,\\[8pt]
0, & A=\emptyset .
\end{cases}
\end{equation}
\end{lemma}

The lemma gives a simple rule:
\begin{itemize}[leftmargin=*]
\item the mass on an admissible value rises with accuracy;
\item the mass on an error value falls;
\item when $\Tell$ is silent, nothing depends on accuracy at all.
\end{itemize}
\Cref{tab:signatures} lists the signs, and \cref{fig:signatures} draws the curves.

The sign difference matters most for ambiguity. A concentrated error $a\notin A$ and the second reading $a\in A$ of an ambiguous cell can have the same frequency at one accuracy level. They move in opposite directions as accuracy grows, which is what separates them.

\begin{table}[t]
\caption{Accuracy-gradient signatures (\cref{lem:signatures}) for two models with $p_1>p_2$ and $\Delta=p_1-p_2$.}
\label{tab:signatures}
\small
\begin{tabular}{@{}lll@{}}
\toprule
Form of $A$ & $D(y)$ for $y\in A$ & $D(y)$ for $y\notin A$ \\
\midrule
$\{x\}$ (no defect if $x=d$; mis-key otherwise) & $\Delta>0$ & $-\Delta\,\tilde\varepsilon(y)\le0$ \\
$\{x,a\}$ (ambiguity) & $\Delta/2>0$ & $-\Delta\,\tilde\varepsilon(y)\le0$ \\
$\emptyset$ (silence) & --- & $0$ for every $y$ \\
\bottomrule
\end{tabular}
\end{table}

\begin{figure}[t]
\centering
\begin{tikzpicture}
\begin{groupplot}[
  group style={group size=4 by 1, horizontal sep=9mm, y descriptions at=edge left},
  width=0.29\textwidth, height=0.25\textwidth,
  xmin=0.3, xmax=1, ymin=0, ymax=1,
  xlabel={accuracy $p$}, ylabel={$P_m(y)$},
  xtick={0.4,0.6,0.8,1}, ytick={0,0.25,0.5,0.75,1},
  tick label style={font=\scriptsize}, label style={font=\scriptsize}, title style={font=\footnotesize},
  every axis plot/.append style={thick, samples=2, domain=0.3:1}]
\nextgroupplot[title={$A=\{x\}$, flat errors}]
  \addplot[black] {x};
  \addplot[black, dashed] {(1-x)/3};
\nextgroupplot[title={$A=\{x,a\}$, flat errors}]
  \addplot[black] {x/2};
  \addplot[blue, densely dotted, very thick] {x/2};
  \addplot[black, dashed] {(1-x)/2};
\nextgroupplot[title={$A=\{x\}$, concentrated error}]
  \addplot[black] {x};
  \addplot[red, densely dashdotted, very thick] {0.84*(1-x)};
  \addplot[black, dashed] {0.08*(1-x)};
\nextgroupplot[title={$A=\emptyset$}]
  \addplot[black, dashed] {0.25};
\end{groupplot}
\end{tikzpicture}
\caption{Accuracy-gradient signatures drawn from \cref{eq:reading} with $|V|=4$. Solid lines show admissible values and dashed lines show error values. In the second panel, the dotted blue line is the second admissible value of an ambiguous cell. In the third panel, the dash-dotted red line is an error value holding $84\%$ of the error mass. At $p\approx0.63$, the blue line in the second panel and the red line in the third have similar heights, but they have opposite slopes. Under silence, no reading depends on accuracy. The curves are analytical, not experimental.}
\label{fig:signatures}
\end{figure}

\subsection{Identifiability}

\begin{theorem}[T1: single value, one model]
\label{thm:t1}
Let $A=\{x\}$, and suppose the plurality condition
\begin{equation}
\label{eq:pc}
p_m>(1-p_m)\max_{y}\tilde\varepsilon(y)\tag{PC$_m$}
\end{equation}
holds. Then $x$ is the unique mode of $P_m$, and $p_m=P_m(x)$.
\end{theorem}

\begin{theorem}[T2: all forms from two models]
\label{thm:t2}
Under (A1)--(A5), with $p_1>p_2$,
\begin{equation}
\label{eq:t2}
A=\{y\in V:D(y)>0\}.
\end{equation}
In particular, $A=\emptyset$ exactly when $D\equiv0$. Under (A4), $\Delta_c>0$ at every cell whenever $\alpha_1\neq\alpha_2$. Hence the forms no defect, mis-key, ambiguity and silence are identified at every cell from any two models of distinct ability.
\end{theorem}

\Cref{thm:t2} needs an accuracy contrast. The flat error pool of \cref{eq:paramprior} gives a second route that needs only one model.

\begin{theorem}[T2$'$: all forms from one model with a known error pool]
\label{thm:t2prime}
Let $\varepsilon=u$ be flat over $V$, with $K=|V|$, and let $p$ be the model's accuracy at the cell.
\begin{itemize}[leftmargin=*]
\item If $A\neq\emptyset$, then $P_m(y)=p/|A|$ for $y\in A$, and $P_m(y)=(1-p)/(K-|A|)$ for $y\notin A$.
\item If $A=\emptyset$, then $P_m$ is flat at $1/K$.
\end{itemize}
Hence $A$ is identified from $P_m$ alone whenever
\begin{equation}
\label{eq:t2prime}
p\neq|A|/K,
\end{equation}
that is, whenever the faithful share differs from what chance would give.
\end{theorem}

\Cref{thm:t1,thm:t2,thm:t2prime} assume that accuracies are known or can be contrasted. The anchors make this possible without labels, as long as most anchors are correct.

\begin{theorem}[T3: model contrasts from anchors]
\label{thm:t3}
At an anchored cell with $A_c=\{d_c\}$, the logit contrast
\begin{equation}
\label{eq:t3}
\ell_c\defeq\logit P_{1,c}(d_c)-\logit P_{2,c}(d_c)=\alpha_1-\alpha_2
\end{equation}
is the same at every non-defective anchored cell. If more than half of the anchored cells are non-defective, then $\median_c\ell_c=\alpha_1-\alpha_2$. The model order and the contrasts are therefore identified, even where \eqref{eq:pc} fails at individual cells: the estimator has a breakdown point of $\tfrac12$. The levels $\alpha_m$ are identified under the location constraint that $\median\beta_{g(c)}=0$ over the non-defective anchored cells.
\end{theorem}

\Cref{thm:t3} is also the reason for one of the EM starts (\cref{sec:approach-model}). When defects are not sparse, it is no longer guaranteed, and the multi-start procedure guards against the resulting poor fixed points.

\subsection{Finite samples}

Population-level identifiability says nothing about how many runs are needed. The next theorem bounds the error of three simple decision rules. Let $n$ be the smaller number of readings per model at the cell, and $\hat P_m$ the empirical reading distributions.

\begin{theorem}[T4: finite-sample bounds]
\label{thm:t4}
Under (A1)--(A5):
\begin{enumerate}[leftmargin=*,label=(\alph*)]
\item \textbf{Two-model contrast.} The rule $\hat A=\{y:\hat D(y)>\Delta/4\}$ errs with probability at most
  \begin{equation}\label{eq:t4a}
  4|V|\exp\bigl(-n\Delta^2/32\bigr).
  \end{equation}
\item \textbf{Plug-in mode.} Under \eqref{eq:pc} with margin $\gamma=p-(1-p)\max_y\tilde\varepsilon(y)$, the mode of $\hat P_m$ differs from $x$ with probability at most
  \begin{equation}\label{eq:t4b}
  2|V|\exp\bigl(-n\gamma^2/2\bigr).
  \end{equation}
\item \textbf{Minimum distance.} Under \cref{thm:t2prime}, let $\gamma_{\min}$ be the smallest sup-norm distance between the predicted reading distributions of two candidate sets. The rule that picks the candidate closest to $\hat P_m$ errs with probability at most
  \begin{equation}\label{eq:t4c}
  2|V|\exp\bigl(-n\gamma_{\min}^2/2\bigr).
  \end{equation}
\end{enumerate}
\end{theorem}

The bounds are conservative: they come from Hoeffding's inequality~\cite{hoeffding1963probability} and a union bound. \specdx does not use these rules directly; it uses the MAP estimator with priors. The simulation study measures its actual power. The bounds still show which quantities matter.
\begin{itemize}[leftmargin=*]
\item \textbf{The accuracy gap and the number of readings.} The error rate falls exponentially in $n\Delta^2$. Pooling lifted cells across scenarios increases $n$.
\item \textbf{Concentrated errors.} An error that holds most of the error mass reduces $\gamma$ and hence the power of the single-model rules. This is why ambiguity, whose signature is weakest, benefits most from several models.
\item \textbf{The size of the candidate set.} The set size $|V|$ enters only logarithmically, so $K_{\max}=6$ costs little.
\end{itemize}

\subsection{Soundness of the logical components}

\begin{proposition}[Sound refutation]
\label{prop:refute}
If refutation returns a state $w$ on which every clause of $M$ evaluates to $\sat$, then $M$ is satisfiable, and $M$ is not a contradiction.
\end{proposition}

\begin{proposition}[No identification without action]
\label{prop:noaction}
Suppose that no acted behaviour touches the rows read by clause $\varphi$. Then every behaviour-based observation has the same distribution when $\varphi$ is a hidden requirement as when $\varphi$ is correct but unreachable by the population. No method that observes only behaviours can distinguish the two cases, and abstaining is the only calibrated output.
\end{proposition}

\Cref{prop:noaction} applies to every behaviour-based baseline in our study as much as to \specdx. It motivates both the Unobservable list of stage 1 and the observability-aware scoring of the evaluation (\cref{sec:eval-design}).

\subsection{Scope of the theory}
\label{sec:theory-scope}

\paragraph{Where the assumptions fail.}
\begin{itemize}[leftmargin=*]
\item \textbf{(A3): errors concentrate or differ by model.} The flat pool is then misspecified, and a concentrated error can look like the second reading of an ambiguous cell. \Cref{lem:signatures} shows that the two still differ in slope. With few models and small accuracy gaps, however, this difference is hard to see. The simulation study includes an arm that violates (A3).
\item \textbf{(A5): a strong model guesses an unstated value.} This only moves a finding from the \emph{unstated} to the \emph{mis-key} subtype. Both subtypes are hidden requirements, so the reported type does not change. The simulation includes an arm that violates (A5).
\end{itemize}

\paragraph{What the theory does not cover.} Contradictions are symbolic and are not covered by \cref{thm:t1,thm:t2,thm:t2prime,thm:t3,thm:t4}; \cref{prop:refute} covers them. Weak checks rest on \cref{thm:t1}: they need only a plurality margin at the cell, not an accuracy contrast.

\paragraph{The advantage depends on the regime.} Simple heuristics can match the anchored cell model in some regimes. Majority vote over models does well when accuracy is moderate and the population is small. Flagging every deviation from the anchor does well on hidden requirements, but it is blind to ambiguity. The evaluation therefore reports a full regime map instead of a single operating point (\cref{sec:results-rq1}).
